% ===== 本次新增：旁切圆几何精确图（原有 assets/tikz 文件完全未改） =====
\begin{center}
\begin{tikzpicture}[x=1.08cm,y=1.08cm,line cap=round,line join=round,>=stealth]
  % 精确数值几何：a=1,b^2=3,c=2，且三边为 6,4,4。
  \pgfmathsetmacro{\hh}{3*sqrt(7)/2}
  \pgfmathsetmacro{\rr}{sqrt(7)}
  \pgfmathsetmacro{\ty}{9*sqrt(7)/8}
  \pgfmathsetmacro{\uy}{7*sqrt(7)/4}
  \coordinate (Fone) at (-2,0);
  \coordinate (Ftwo) at (2,0);
  \coordinate (P) at (2.5,\hh);
  \coordinate (T) at (2.375,\ty);
  \coordinate (U) at (3.25,\uy);
  \coordinate (V) at (5,0);
  \coordinate (I) at (5,\rr);
  % 双曲线准确为 x^2-y^2/3=1，点 P 恰落在右支。
  \draw[blue!55,line width=0.8pt,domain=-0.55:4.85,samples=100,smooth,variable=\z]
     plot ({sqrt(1+(\z)^2/3)}, {\z});
  \draw[blue!55,line width=0.8pt,domain=-0.55:4.85,samples=100,smooth,variable=\z]
     plot ({-sqrt(1+(\z)^2/3)}, {\z});
  % 坐标轴：水平直线即焦点连线的支撑直线。
  \draw[black!45,->] (-2.50,0)--(8.0,0) node[right] {$x$};
  \draw[black!45,->] (0,-0.5)--(0,5.55) node[above] {$y$};
  \fill[red!8] (Fone)--(Ftwo)--(P)--cycle;
  % 旁切圆：I=(5,sqrt(7)), radius=sqrt(7)。
  \fill[green!5] (I) circle (\rr);
  \draw[green!55!black,line width=0.95pt] (I) circle (\rr);
  % 三条边及两条切线所在边的延长部分。
  \draw[red!80!black,line width=1.2pt] (Fone)--(P)--(Ftwo)--cycle;
  \draw[red!75!black,dashed,line width=0.9pt] (P)--(U);
  \draw[red!75!black,dashed,line width=0.9pt] (Ftwo)--(V);
  % 三条半径，全部垂直于切线。
  \draw[purple!75,densely dashed] (I)--(T) (I)--(U) (I)--(V);
  % 从 P 到圆心：顶角的外角平分线。
  \draw[green!55!black,densely dotted,line width=0.9pt] (P)--(I);
  % 三个直角标记：T处向P（PF2）、U处向P（PF1）、V处向F2（x轴）。
  \draw[purple!75,line width=0.65pt] (T) ++(0.025,{3*sqrt(7)/40})
       -- ++({3*sqrt(7)/40},-0.025) -- ++(-0.025,{-3*sqrt(7)/40});
  \draw[purple!75,line width=0.65pt] (U) ++(-0.15,{-sqrt(7)/20})
       -- ++({sqrt(7)/20},-0.15) -- ++(0.15,{sqrt(7)/20});
  \draw[purple!75,line width=0.65pt] (4.80,0)--(4.80,0.20)--(5,0.20);
  % theta 为 P 点三角形内部夹角。
  \draw[red!80!black,thin] (P) ++(-97.1808:0.49) arc (-97.1808:-138.5904:0.49);
  \node[font=\small] at (2.02,3.32) {$\theta$};
  \foreach \q in {Fone,Ftwo,P,T,U,V} \fill[red!80!black] (\q) circle (1.65pt);
  \fill[green!40!black] (I) circle (1.75pt);
  \node[below left,font=\small] at (Fone) {$F_1$};
  \node[below left,font=\small] at (Ftwo) {$F_2$};
  \node[above left,font=\small] at (P) {$P$};
  \node[left,font=\small] at (T) {$T$};
  \node[above right,font=\small] at (U) {$U$};
  \node[below,font=\small] at (V) {$V$};
  \node[above right,font=\small] at (I) {$I$};
  % 数值为 PU=PT=c-a=1, F2T=F2V=3，避免图面拥挤，长度在图下证明中标注。
  \node[font=\scriptsize,purple!75,right] at (5.02,1.35) {$r$};
  \node[font=\small,green!40!black,align=center] at (6.05,4.05) {旁切圆\\(对应 $F_1$)};
\end{tikzpicture}
\end{center}
% ===== 精确图结束 =====
